% Gerado a partir de Unidade_1/Resolucoes_e_Pratica/Lista_Treino_Facil.md
\chapter{Lista de Treino - Nível Fácil (15 Questões)}

\textbf{Disciplina:} Sinais e Sistemas (TEC 412)\\
\textbf{Tópicos:} Paridade, Energia/Potência, Memória, Causalidade, Linearidade, Invariância e Estabilidade.

\section{Tópico 1: Sinais Pares e Ímpares}\label{tuxf3pico-1-sinais-pares-e-uxedmpares}

\textbf{Questão 1:} Dado o polinômio \(x(t) = 1 + t + t^2 + t^3\), determine analiticamente a sua parte par \(x_p(t)\) e sua parte ímpar \(x_i(t)\).

\begin{quote}
\textbf{Resolução:}\\
Parte Par: \(x_p(t) = \frac{x(t) + x(-t)}{2}\)\\
\(x_p(t) = \frac{(1 + t + t^2 + t^3) + (1 + (-t) + (-t)^2 + (-t)^3)}{2}\)\\
\(x_p(t) = \frac{1 + t + t^2 + t^3 + 1 - t + t^2 - t^3}{2}\)\\
\(x_p(t) = \frac{2 + 2t^2}{2} = 1 + t^2\).\\
Parte Ímpar: \(x_i(t) = \frac{x(t) - x(-t)}{2}\)\\
\(x_i(t) = \frac{(1 + t + t^2 + t^3) - (1 - t + t^2 - t^3)}{2}\)\\
\(x_i(t) = \frac{2t + 2t^3}{2} = t + t^3\).
\end{quote}

\textbf{Questão 2:} Analise graficamente e comprove algebricamente se o sinal retangular que vai de \(t=-2\) até \(t=2\) com amplitude constante \(1\) é par ou ímpar. \(x(t) = u(t+2) - u(t-2)\).

\begin{quote}
\textbf{Resolução:}\\
Um sinal contínuo igual a 1 de \(-2\) a \(2\) apresenta simetria perfeita em relação ao eixo vertical \(Y\).\\
Comprovando algebricamente:\\
\(x(-t) = u(-t+2) - u(-t-2)\).\\
Como a função retangular é estritamente simétrica em torno do zero, inverter o tempo não muda os seus limites físicos no espaço. A amplitude em \(t=1\) é \(1\), e em \(t=-1\) também é \(1\). Logo, \(x(-t) = x(t)\).\\
\textbf{Gabarito: Sinal Par.}
\end{quote}

\section{Tópico 2: Energia e Potência}\label{tuxf3pico-2-energia-e-potuxeancia}

\textbf{Questão 3:} Calcule a energia total do pulso retangular \(x(t) = 3\) definido apenas no intervalo \(0 \le t \le 4\), sendo zero no resto do tempo.

\begin{quote}
\textbf{Resolução:}\\
\(E = \int_{-\infty}^{\infty} |x(t)|^2 dt = \int_{0}^{4} (3)^2 dt\)\\
\(E = \int_{0}^{4} 9 dt = [9t]_0^4 = 9(4) - 9(0) = 36 \text{ Joules}\).\\
\textbf{Gabarito: Energia = 36 Joules. Como é finita, é um sinal de Energia.}
\end{quote}

\textbf{Questão 4:} O sinal contínuo constante \(x(t) = 5\) estendido de \(-\infty\) a \(\infty\) é um sinal de energia ou potência? Determine o valor.

\begin{quote}
\textbf{Resolução:}\\
A integral de energia de uma constante ao longo do tempo infinito será infinita. Portanto, testamos a potência.\\
\(P = \lim_{T \to \infty} \frac{1}{2T} \int_{-T}^{T} (5)^2 dt = \lim_{T \to \infty} \frac{1}{2T} \int_{-T}^{T} 25 dt\)\\
\(P = \lim_{T \to \infty} \frac{1}{2T} [25t]_{-T}^{T} = \lim_{T \to \infty} \frac{1}{2T} (25T - (-25T)) = \lim_{T \to \infty} \frac{50T}{2T} = 25\).\\
\textbf{Gabarito: Sinal de Potência, P = 25 Watts.}
\end{quote}

\section{Tópico 3: Memória}\label{tuxf3pico-3-memuxf3ria}

\textbf{Questão 5:} O sistema \(y(t) = 4x(t) + 2\) possui memória?

\begin{quote}
\textbf{Resolução:}\\
Avaliando num instante \(t_0\), temos \(y(t_0) = 4x(t_0) + 2\). A saída exige unicamente a entrada no mesmo instante \(t_0\). Não há integrais, derivadas ou deslocamentos.\\
\textbf{Gabarito: Sistema SEM Memória (Estático).}
\end{quote}

\textbf{Questão 6:} O sistema contínuo \(y(t) = x(t^2)\) possui memória?

\begin{quote}
\textbf{Resolução:}\\
Vamos testar um instante de tempo. Por exemplo, em \(t = 2\):\\
\(y(2) = x(2^2) = x(4)\).\\
Para gerar a saída no segundo 2, o sistema exige saber o valor do sinal no segundo 4. Como precisou ler um instante diferente do momento atual, necessita de armazenamento.\\
\textbf{Gabarito: Sistema COM Memória.}
\end{quote}

\section{Tópico 4: Causalidade}\label{tuxf3pico-4-causalidade}

\textbf{Questão 7:} O sistema \(y(t) = x(t-5)\) é causal?

\begin{quote}
\textbf{Resolução:}\\
A saída num instante \(t_0\) é \(y(t_0) = x(t_0 - 5)\).\\
Como \(t_0 - 5\) é um valor de tempo SEMPRE anterior a \(t_0\), o sistema só depende de informações que já ocorreram (passado).\\
\textbf{Gabarito: Sistema Causal.}
\end{quote}

\textbf{Questão 8:} O sistema de tempo discreto \(y[n] = x[n+1] + x[n]\) respeita a causalidade física?

\begin{quote}
\textbf{Resolução:}\\
Para \(n=0\): \(y[0] = x[1] + x[0]\).\\
O sistema precisa ler a amostra \(x[1]\) para gerar a saída atual no instante 0. Como \(x[1]\) está no futuro em relação ao tempo presente \(n=0\), isso não é viável em tempo real.\\
\textbf{Gabarito: Sistema NÃO-Causal.}
\end{quote}

\section{Tópico 5: Linearidade}\label{tuxf3pico-5-linearidade}

\textbf{Questão 9:} Prove analiticamente se o sistema amplificador \(y(t) = 5x(t)\) é linear.

\begin{quote}
\textbf{Resolução:}\\
Entrada combinada: \(x_3(t) = ax_1(t) + bx_2(t)\).\\
Saída resultante: \(y_3(t) = 5(ax_1(t) + bx_2(t)) = 5ax_1(t) + 5bx_2(t)\).\\
Rearranjando: \(y_3(t) = a(5x_1(t)) + b(5x_2(t)) = a \cdot y_1(t) + b \cdot y_2(t)\).\\
A igualdade se manteve perfeita.\\
\textbf{Gabarito: Sistema Linear.}
\end{quote}

\textbf{Questão 10:} O sistema \(y(t) = x(t) + 1\) é linear?

\begin{quote}
\textbf{Resolução:}\\
Um requisito básico da linearidade (homogeneidade) é a propriedade da aditividade na origem: \(x(t) = 0 \implies y(t) = 0\).\\
Inserindo uma entrada nula: \(y(t) = 0 + 1 = 1\).\\
Como a saída para o repouso não foi nula, a propriedade de superposição fatalmente falhará.\\
\textbf{Gabarito: Sistema NÃO-Linear.}
\end{quote}

\section{Tópico 6: Invariância no Tempo}\label{tuxf3pico-6-invariuxe2ncia-no-tempo}

\textbf{Questão 11:} O atenuador ideal \(y(t) = \frac{1}{2}x(t)\) é invariante no tempo?

\begin{quote}
\textbf{Resolução:}\\
Passo 1 (atrasa entrada): \(y_1(t) = \frac{1}{2}x(t-t_0)\).\\
Passo 2 (atrasa variável global): A única variável global está no argumento do \(x\). Logo, \(y_2(t) = \frac{1}{2}x(t-t_0)\).\\
Como \(y_1(t) = y_2(t)\), o sistema não muda suas características com o passar do relógio.\\
\textbf{Gabarito: Sistema Invariante no Tempo.}
\end{quote}

\textbf{Questão 12:} O sistema oscilador \(y(t) = x(t) \cdot \cos(3t)\) é invariante?

\begin{quote}
\textbf{Resolução:}\\
Atrasando a entrada: \(y_1(t) = x(t-t_0) \cdot \cos(3t)\).\\
Atrasando a referência temporal do sistema: substituímos \(t\) por \(t-t_0\):\\
\(y_2(t) = x(t-t_0) \cdot \cos(3(t-t_0))\).\\
Vemos claramente que \(y_1 \neq y_2\) porque a fase do cosseno mudou no segundo caso.\\
\textbf{Gabarito: Sistema Variante no Tempo.}
\end{quote}

\section{Tópico 7: Estabilidade BIBO}\label{tuxf3pico-7-estabilidade-bibo}

\textbf{Questão 13:} O sistema trigonométrico \(y(t) = \cos(x(t))\) é BIBO estável?

\begin{quote}
\textbf{Resolução:}\\
Assuma entrada limitada \(|x(t)| \le B_x\).\\
O sinal \(y(t)\) é o cosseno de alguma coisa. Sabemos que a função cosseno, matematicamente, sempre retorna valores restritos entre \(-1\) e \(+1\), independentemente do tamanho do ângulo interno.\\
Portanto, a amplitude da saída \(|y(t)| \le 1\).\\
Como 1 é sempre um número finito, o sistema jamais explodirá.\\
\textbf{Gabarito: Sistema Estável.}
\end{quote}

\textbf{Questão 14:} O sistema de rampa inversa \(y(t) = \frac{x(t)}{t}\) (com \(t \ge 0\)) é estável?

\begin{quote}
\textbf{Resolução:}\\
Vamos testar uma entrada constante e limitada, por exemplo \(x(t) = 1\) (degrau).\\
A saída seria \(y(t) = \frac{1}{t}\).\\
Quando avaliamos o sistema nas vizinhanças de \(t=0\), o limite de \(\frac{1}{t}\) para \(t \to 0\) vai para o infinito (\(\infty\)).\\
Tivemos uma entrada que nunca passou de 1, mas a saída explodiu em \(t=0\).\\
\textbf{Gabarito: Sistema Instável.}
\end{quote}

\section{Tópico 8: Mistura de Conceitos (Sistemas Completos)}\label{tuxf3pico-8-mistura-de-conceitos-sistemas-completos}

\textbf{Questão 15:} Dado o sistema de tempo discreto \(y[n] = n \cdot x[n]\). Classifique-o quanto a Memória, Causalidade, Linearidade e Invariância.

\begin{quote}
\textbf{Resolução e Gabarito:}

\begin{enumerate}
\def\labelenumi{\arabic{enumi}.}
\tightlist
\item
  \textbf{Memória:} Para o instante \(n\), só pede \(x[n]\). Logo, \textbf{Sem Memória}.
\item
  \textbf{Causalidade:} Como é sem memória (só pede o presente), é automaticamente \textbf{Causal}.
\item
  \textbf{Linearidade:} \(y_3[n] = n \cdot (ax_1[n] + bx_2[n]) = a(n \cdot x_1[n]) + b(n \cdot x_2[n]) = a \cdot y_1[n] + b \cdot y_2[n]\). \textbf{Linear}.
\item
  \textbf{Invariância:} Atrasando a entrada: \(y_1[n] = n \cdot x[n-n_0]\). Atrasando o sistema geral: \(y_2[n] = (n-n_0) \cdot x[n-n_0]\). Não são iguais. Logo, é \textbf{Variante no Tempo}.
\end{enumerate}
\end{quote}
